\subsection{Purpose}

This section defines requirements for protecting facilities, equipment, systems, media, and information from unauthorized physical access, theft, damage, tampering, environmental hazards, and service interruption.

\subsection{Scope}

These requirements apply to all organizational offices, work areas, server rooms, equipment rooms, storage areas, backup locations, and other facilities where organizational systems or information are located.

They also apply to organizational equipment used at homes, remote offices, customer locations, public locations, or other facilities.

\subsection{Physical Access}

Physical access to areas containing systems or sensitive information \must{} be limited to authorized individuals.

Access controls \should{} include appropriate measures such as:
\begin{itemize}
    \item Keys, access cards, badges, or electronic locks
    \item Visitor sign-in and escort procedures
    \item Physical barriers and locked cabinets
    \item Security cameras or alarm systems where appropriate
    \item Access logs or other records of entry
\end{itemize}

Access rights \must{} be granted according to job responsibilities and business need.

Physical access \must{} be removed promptly when a person leaves the organization, changes roles, or no longer requires access.

Keys, access cards, badges, and other physical access devices \must{} be returned, disabled, or replaced when they are lost, stolen, or no longer required.

\subsection{Visitors and Contractors}

Visitors and contractors \must{} be authorized before entering restricted areas.

Visitors \must{}:
\begin{itemize}
    \item Sign in when required
    \item Wear visitor identification when appropriate
    \item Remain accompanied or monitored in restricted areas
    \item Access only the areas necessary for their approved purpose
    \item Return temporary access devices before leaving
\end{itemize}

Visitors and contractors \mustnot{} be left unsupervised in areas containing sensitive systems or information unless their access has been specifically approved.

\subsection{Restricted Areas}

Server rooms, network equipment rooms, backup storage areas, and other restricted locations \must{} be secured against unauthorized access.

Restricted areas \must{}:
\begin{itemize}
    \item Remain locked when unattended
    \item Be protected from unauthorized equipment removal
    \item Be kept free of unnecessary combustible materials
    \item Have access limited to authorized personnel
    \item Be monitored for environmental or equipment risks where appropriate
\end{itemize}

Equipment, cables, network ports, and storage media in restricted areas \must{} be protected against tampering or unauthorized connection.

\subsection{Equipment Protection}

Organizational equipment \must{} be protected against theft, unauthorized use, damage, and environmental hazards.

Workstations and laptops containing organizational information \should{} be secured using appropriate physical measures, including:
\begin{itemize}
    \item Locking devices or secure storage
    \item Screen locks
    \item Cable locks where appropriate
    \item Secure transportation
    \item Protection from unauthorized viewing
\end{itemize}

Devices \must{} not be left unattended in vehicles or other locations where theft or unauthorized access is reasonably likely.

Equipment used in public or shared locations \must{} be positioned to prevent unauthorized viewing or access whenever practical.

\subsection{Clear Desk and Clear Screen}

Users \should{} keep desks and work areas clear of Confidential and Restricted information when the information is not actively being used.

Confidential and Restricted documents \must{} be stored in locked cabinets, secure rooms, or other approved locations when unattended.

Users \must{} lock their screens when leaving devices unattended.

Printed sensitive information \mustnot{} be left in printers, meeting rooms, reception areas, or other publicly accessible locations.

\subsection{Environmental Protection}

Critical systems and equipment \must{} be protected from environmental conditions that could cause loss, damage, or interruption.

Where appropriate, facilities \should{} provide protection against:
\begin{itemize}
    \item Fire, smoke, and water
    \item Excessive heat or humidity
    \item Power interruption or electrical damage
    \item Unauthorized removal
    \item Dust and other contaminants
    \item Natural disasters and other facility hazards
\end{itemize}

Critical equipment \should{} use appropriate power protection, such as surge protection, backup power, or an uninterruptible power supply.

Environmental alarms and facility failures \must{} be reported and investigated promptly when they could affect systems or information.

\subsection{Removable Media}

Removable media \must{} be approved before it is connected to organizational systems when the media is not organization-owned or trusted.

Removable media containing Confidential or Restricted information \must{} be encrypted when technically feasible.

Removable media \must{} be stored securely when not in use and \must{} be securely disposed of when no longer required.

Unknown or suspicious removable media \mustnot{} be connected to organizational systems.

\subsection{Remote and Home Work Areas}

Personnel working remotely \must{} take reasonable steps to prevent unauthorized access, disclosure, theft, or damage.

Remote work areas \should{}:
\begin{itemize}
    \item Prevent unauthorized people from viewing screens or documents
    \item Keep organizational devices physically secure
    \item Use private and trusted workspaces when handling sensitive information
    \item Prevent family members, guests, or other unauthorized individuals from using organizational devices
    \item Secure paper records and removable media
\end{itemize}

Organizational equipment \must{} be transported and stored in a manner that reduces the risk of loss, theft, or damage.

\subsection{Equipment Removal and Disposal}

Equipment \must{} be authorized before it is removed from organizational facilities.

Before equipment is transferred, returned, recycled, or disposed of, storage media \must{} be securely erased or physically destroyed according to the data classification and disposal requirements.

Physical access devices, labels, documents, storage media, and other items containing organizational information \must{} be returned or securely destroyed when no longer required.

\subsection{Physical Security Incidents}

Lost keys, badges, devices, media, documents, or other physical assets \must{} be reported promptly.

Suspected unauthorized entry, theft, tampering, environmental damage, or physical security control failure \must{} be reported according to the incident response procedures.

The organization \should{} document investigations and corrective actions for significant physical security incidents.

\subsection{Review}

Physical security controls \should{} be reviewed at least annually and whenever facilities, equipment, personnel, system requirements, or threats change.

Access lists for restricted areas \must{} be reviewed periodically to verify that access remains appropriate.